\begin{frame}{그림으로: 오토인코더 vs VAE의 잠재 공간}
  \begin{center}
  \includegraphics[height=5.6cm]{gen_ae_vs_vae_latent.png}
  \end{center}
  \vskip0.3em
  \small
  일반 오토인코더(왼쪽)는 학습점이 여기저기 뭉쳐 있고 그 사이 \textbf{빈
  공간}(빨간 $\times$)에서 뭐가 나올지 보장이 없다. VAE(오른쪽)는 모든
  분포를 $\mathcal{N}(0,1)$ 쪽으로 당겨 공간이 \textbf{매끄럽고
  빈틈없이} 채워진다 --- 같은 위치(빨간 별)에서도 VAE 쪽만 항상
  데이터 근처.
\end{frame}
